
% ============================================================
\section{Research Methodology}
% ============================================================

\subsection{Design Approach}

The project follows an iterative engineering design cycle:
requirements were derived from the deficiencies in
Section~I-B, an architecture was specified, each module was
implemented and exercised through automated tests, and evaluation
fed back into hardening. Requirements are stated as testable
claims (e.g., ``the overlay never uses \texttt{win.hide()}''),
so that every documented behaviour is backed by an executable
check rather than by inspection alone.

\subsection{Technology Selection}

Table~\ref{tab:stack} lists the selected technologies and the
criteria behind each choice.

\begin{table}[t]
\centering\small
\caption{Technology selection and criteria}\label{tab:stack}
\begin{tabular}{@{}p{0.26\columnwidth}p{0.44\columnwidth}p{0.22\columnwidth}@{}}
\toprule
\textbf{Layer} & \textbf{Technology} & \textbf{Criterion} \\
\midrule
Overlay UI & Electron transparent window &
Native click-through, CSS control \\
Remote channel & WhatsApp Web bridge &
Phone-first, no new app to install \\
Inference (hosted) & OpenCode session API &
Existing account, free-tier models \\
Inference (cloud) & OpenAI-compatible endpoints &
One client, many providers \\
Inference (local) & Ollama (localhost) &
Offline availability \\
Scheduling & Windows Task Scheduler &
Survives logout/reboot \\
\bottomrule
\end{tabular}
\end{table}

The OpenAI-compatible criterion deserves emphasis: a single
HTTP client speaking the de-facto \texttt{/v1/chat/completions}
dialect covers OpenAI, OpenRouter, DeepSeek, Groq and any
self-hosted gateway, which is what makes ``works with all
providers'' an architectural property rather than a per-vendor
integration project.

\subsection{Evaluation Corpus}

Because the contribution is a system rather than a statistical
model, the evaluation corpus is the project's automated test
suite together with controlled failure injections:

\begin{itemize}
  \item \textbf{Test suite:} 27 modules under
  \texttt{scripts/test-*.js}, each driving one subsystem through
  its public interface (queue discipline, permission gating,
  digest generation, news parsing, model probing, repository
  watching, media control, self-audit, screenshot gallery, and
  overlay mode transitions).
  \item \textbf{Failure injections:} (i)~provider quota
  exhaustion --- the session API rejects with a credit error;
  (ii)~service unreachable --- \texttt{ECONNREFUSED} on the local
  service port; (iii)~network loss --- \texttt{ENOTFOUND} on
  external hosts. Each injection asserts that a reply is still
  produced by the next tier of the chain.
  \item \textbf{Operational telemetry:} the append-only debug log
  records every prompt, fallback decision, queue transition and
  watchdog action, providing the qualitative evidence used in
  Section~V.
\end{itemize}

\subsection{Exclusion Criteria}

Behaviours not demonstrated by a passing test are excluded from
reported results. Non-deterministic demonstrations (manual
click-throughs, one-off screenshots) are treated as illustrative
only. This mirrors the exclusion of non-peer-reviewed sources in
a conventional literature review: only evidence that can be
re-executed is counted.

% ============================================================
\section{Analysis of the System}
% ============================================================

The analysis follows two thematic parts: the \emph{performance
metrics} the system is built to optimize
(Section~IV-A) and the \emph{inference tiers and modules} that
realize them (Section~IV-B).

\subsection{Performance Metrics}

Goals are grouped into three families, mirroring the three
prediction targets of classic educational-data studies
(completion, grades, at-risk early warning) but instantiated as
system-level service metrics.

\subsubsection{Availability \& Continuity (completion analog)}
\begin{itemize}
  \item \textbf{Reply availability:} fraction of user prompts
  that receive a substantive reply rather than an error. The
  design target is 100\% while at least one tier is reachable,
  and a non-error deterministic answer when none is.
  \item \textbf{Offline-path success rate:} under injected
  network loss, share of prompts answered by the local or
  deterministic tier.
  \item \textbf{Watchdog recovery time:} elapsed time between a
  host-process death and the scheduled relaunch; the watchdog
  runs every minute with missed-run recovery enabled, bounding
  recovery to one scheduling interval after wake.
\end{itemize}

\subsubsection{Responsiveness (grades analog)}
\begin{itemize}
  \item \textbf{Reply latency:} prompt-to-first-reply time per
  tier (hosted $<$ cloud $<$ local in typical measurements).
  \item \textbf{Queue throughput:} completions per minute under
  sustained phone traffic; the strict FIFO with capacity
  $Q_{max}=8$ bounds memory and gives the operator an explicit
  back-pressure signal (``queue full'') instead of silent
  unbounded growth.
  \item \textbf{Digest freshness:} age of the newest persisted
  news item at delivery time.
\end{itemize}

\subsubsection{Safety \& Reliability (at-risk analog)}
\begin{itemize}
  \item \textbf{Permission-gating coverage:} share of
  consequential actions (shell commands, file sends, provider
  changes) that pass through an \emph{allow/always/reject}
  prompt before execution; \emph{always} whitelists are recorded
  so the audit trail shows why an action ran without a prompt.
  \item \textbf{Execution exclusivity:} invariant that at most
  one action executes at a time --- verified by the queue test.
  \item \textbf{Test pass rate:} aggregate result of the
  27-module suite, used as the primary regression metric.
\end{itemize}

\subsection{Inference Tiers \& System Modules}

\subsubsection{Provider-Agnostic Fallback Chain}

The inference layer is a four-tier chain
(Fig.~\ref{fig:chain}); on failure at tier $n$ the router
advances to $n{+}1$ and records the reason in the debug log.

\begin{figure}[t]
\centering
\begin{tikzpicture}[
  node distance=6mm and 2mm,
  tier/.style={draw, rounded corners=3pt, align=center,
    minimum width=0.92\columnwidth, inner sep=5pt,
    fill=blue!6},
  fail/.style={-{Stealth[length=2mm]}, dashed, gray!70!black,
    font=\scriptsize}
]
\node[tier] (t1) {\textbf{Tier 1 --- Hosted session}\\
OpenCode session API · free models
(\texttt{mimo}, \texttt{deepseek}, \texttt{hy4}) with in-chain
rotation};
\node[tier, below=of t1] (t2) {\textbf{Tier 2 --- OpenAI-compatible cloud}\\
OpenRouter / DeepSeek / Groq / any
\texttt{baseURL} + key};
\node[tier, below=of t2] (t3) {\textbf{Tier 3 --- Local runtime}\\
Ollama at \texttt{127.0.0.1:11434}};
\node[tier, below=of t3] (t4) {\textbf{Tier 4 --- Deterministic offline}\\
Status, cached weather/news, queue state};
\draw[fail] (t1) -- node[right]{4xx/credit} (t2);
\draw[fail] (t2) -- node[right]{5xx/timeout} (t3);
\draw[fail] (t3) -- node[right]{no daemon} (t4);
\end{tikzpicture}
\caption{Four-tier provider-agnostic fallback chain. A failure
at any tier advances the router and is logged with its
reason.}\label{fig:chain}
\end{figure}

Tier~1 rotates across hosted free models when a quota error is
detected, preserving the historical behaviour of the project;
Tier~2 is provider-generic by construction (one HTTP dialect);
Tier~3 restores full generativity without any network path;
Tier~4 never fails because it performs no I/O beyond reading
state the system already holds. The chain is evaluated in
Section~V.

\subsubsection{Feature Modules}

Table~\ref{tab:modules} maps each service module to its
responsibility and its governing test.

\begin{table}[t]
\centering\small
\caption{Service modules and governing tests}\label{tab:modules}
\begin{tabular}{@{}p{0.24\columnwidth}p{0.46\columnwidth}p{0.2\columnwidth}@{}}
\toprule
\textbf{Module} & \textbf{Responsibility} & \textbf{Test} \\
\midrule
Overlay & Pill/panel, hover expand, pin, dim modes & \texttt{ctrl} \\
Bridge & FIFO queue, commands, permission prompts & \texttt{queue}, \texttt{permission} \\
Assistant & Session, fallback chain, replies & \texttt{continue}, \texttt{actions} \\
News & HN + RSS $\rightarrow$ digest packs & \texttt{news}, \texttt{digest} \\
Weather & Forecast + offline cache & \texttt{weather} \\
GitHub watch & Issues/PRs/CI glow alerts & \texttt{ghwatch} \\
Model probe & Health of each provider tier & \texttt{modelprobe} \\
Media & Transport keys, now-playing pill & (media keys) \\
Self-audit & Log tail anomaly reader & \texttt{selfaudit} \\
Watchdog & Scheduled relaunch + battery policy & registration script \\
\bottomrule
\end{tabular}
\end{table}

The human-in-the-loop path is deliberately narrow: a phone
message enters the bounded queue, at most one action executes,
consequential actions surface a permission prompt, and the
resolution (including \emph{always} rules) is appended to the
audit log. Drag-and-drop file sends and screenshot forwards
reuse the same gated channel rather than opening a second,
ungated path.
